\chapter{Let's Talk about Your Childhood}

% ================= PERCAKAPAN 1-5 =================
\begin{convobox}{1}
    \noindent
    \jpkana{こんにちは。}{Konnichiwa.}

    \vspace{1.5em}
    \noindent{\Large \color{articolor} \textbf{Arti:} Halo (selamat siang).}
\end{convobox}

\begin{convobox}{2}
    \noindent
    \jpword{ひさ}{久}{Hisa}\jpkana{しぶりに}{shiburi ni} 
    \jpword{こうえん}{公園}{kouen} \jpkana{に}{ni} 
    \jpword{き}{来}{ki}\jpkana{ました。}{mashita.}

    \vspace{1.5em}
    \noindent{\Large \color{articolor} \textbf{Arti:} Sudah lama tidak datang ke taman.}
\end{convobox}

\begin{convobox}{3}
    \noindent
    \jpword{きょう}{今日}{Kyou} \jpkana{は}{wa} 
    \jpkana{ここで、}{koko de,} 
    \jpword{こ}{子}{ko}\jpkana{どものころの}{domo no koro no} 
    \jpword{はなし}{話}{hanashi} \jpkana{をしましょう。}{o shimashou.}

    \vspace{1.5em}
    \noindent{\Large \color{articolor} \textbf{Arti:} Hari ini di sini, mari kita bicarakan tentang masa kecil.}
\end{convobox}

\begin{convobox}{4}
    \noindent
    \jpword{こ}{子}{ko}\jpkana{どものころ、}{domo no koro,} 
    \jpkana{どこに}{doko ni} 
    \jpword{す}{住}{su}\jpkana{んでいましたか？}{nde imashita ka?}

    \vspace{1.5em}
    \noindent{\Large \color{articolor} \textbf{Arti:} Sewaktu kecil, kamu tinggal di mana?}
\end{convobox}

\begin{convobox}{5}
    \noindent
    \jpword{おおさか}{大阪}{Oosaka} \jpkana{に}{ni} 
    \jpword{す}{住}{su}\jpkana{んでいました。}{nde imashita.}

    \vspace{1.5em}
    \noindent{\Large \color{articolor} \textbf{Arti:} Dulu tinggal di Osaka.}
\end{convobox}

% --- LESSON BUFFER 1 ---
\begin{lessonbox}{Catatan Belajar: Tata Bahasa "Masa Lalu" (Bentuk Lampau)}
    Hai adik-adik! Pernahkah kamu bercerita tentang kejadian waktu kamu masih kecil? Di bahasa Jepang, kita pakai kata ajaib \textbf{子どものころ (Kodomo no koro)} yang artinya "Sewaktu masih kecil".\\
    
    Nah, karena ini cerita zaman dulu, perhatikan deh kalimat nomor 4 dan 5. Ada kata \textbf{〜ていました (〜te imashita)}. Ini adalah grammar untuk menyatakan \textbf{keadaan yang berlangsung di masa lalu}.
    \begin{itemize}
        \item \textit{Sunde imasu} (Sekarang sedang tinggal) $\rightarrow$ \textbf{Sunde imashita} (Dulu tinggal).
    \end{itemize}
    Coba ingat-ingat, waktu kecil kamu \textit{sunde imashita} di mana?
\end{lessonbox}

% ================= PERCAKAPAN 6-10 =================
\begin{convobox}{6}
    \noindent
    \jpword{なに}{何}{Nani} \jpkana{を}{o} 
    \jpkana{して}{shite} 
    \jpword{あそ}{遊}{aso}\jpkana{ぶのが}{bu no ga} 
    \jpword{す}{好}{su}\jpkana{きでしたか？}{ki deshita ka?}

    \vspace{1.5em}
    \noindent{\Large \color{articolor} \textbf{Arti:} Dulu suka bermain apa?}
\end{convobox}

\begin{convobox}{7}
    \noindent
    \jpword{こうえん}{公園}{Kouen} \jpkana{で}{de} 
    \jpkana{サッカー}{sakkaa} \jpkana{を}{o} 
    \jpkana{して}{shite} 
    \jpword{あそ}{遊}{aso}\jpkana{ぶのが}{bu no ga} 
    \jpword{す}{好}{su}\jpkana{きでした。}{ki deshita.}

    \vspace{1.5em}
    \noindent{\Large \color{articolor} \textbf{Arti:} Dulu suka bermain sepak bola di taman.}
\end{convobox}

\begin{convobox}{8}
    \noindent
    \jpword{そと}{外}{Soto} \jpkana{で}{de} 
    \jpword{あそ}{遊}{aso}\jpkana{ぶのと、}{bu no to,} 
    \jpword{いえ}{家}{ie} \jpkana{の}{no} 
    \jpword{なか}{中}{naka} \jpkana{で}{de} 
    \jpword{あそ}{遊}{aso}\jpkana{ぶの、}{bu no,}

    \vspace{1.5em}
    \noindent{\Large \color{articolor} \textbf{Arti:} Antara bermain di luar dan bermain di dalam rumah,}
\end{convobox}

\begin{convobox}{9}
    \noindent
    \jpkana{どっちの}{Dotchi no} 
    \jpword{ほう}{方}{hou} \jpkana{が}{ga} 
    \jpword{す}{好}{su}\jpkana{きでしたか？}{ki deshita ka?}

    \vspace{1.5em}
    \noindent{\Large \color{articolor} \textbf{Arti:} kamu lebih suka yang mana?}
\end{convobox}

\begin{convobox}{10}
    \noindent
    \jpword{そと}{外}{Soto} \jpkana{で}{de} 
    \jpword{あそ}{遊}{aso}\jpkana{ぶ}{bu} 
    \jpword{ほう}{方}{hou} \jpkana{が}{ga} 
    \jpword{す}{好}{su}\jpkana{きでした。}{ki deshita.}

    \vspace{1.5em}
    \noindent{\Large \color{articolor} \textbf{Arti:} Dulu aku lebih suka bermain di luar.}
\end{convobox}

% --- LESSON BUFFER 2 ---
\begin{lessonbox}{Catatan Belajar: Tata Bahasa Memilih 2 Pilihan (A vs B)}
    Bagaimana cara bertanya "Kamu lebih suka A atau B?" dalam bahasa Jepang? Kita bisa gunakan rumus grammar yang super keren ini:\\
    
    \textbf{[Pilihan A] to, [Pilihan B] to, dotchi no hou ga [kata sifat] desu ka?}\\
    
    Contoh di atas: Bermain di luar (\textit{Soto de asobu no}) VS Bermain di dalam (\textit{Ie no naka de asobu no}). Kata \textbf{どっち (Dotchi)} artinya "Yang mana". Untuk menjawabnya, sebutkan pilihan yang kamu suka dan tambahkan \textbf{〜ほうが (〜hou ga)}. Wah, kamu semakin jago merangkai kalimat panjang!
\end{lessonbox}

% ================= PERCAKAPAN 11-18 =================
\begin{convobox}{11}
    \noindent
    \jpword{こ}{子}{ko}\jpkana{どものころ、}{domo no koro,} 
    \jpkana{ペット}{petto} \jpkana{を}{o} 
    \jpword{か}{飼}{ka}\jpkana{っていましたか？}{tte imashita ka?}

    \vspace{1.5em}
    \noindent{\Large \color{articolor} \textbf{Arti:} Waktu kecil, apakah kamu memelihara hewan peliharaan?}
\end{convobox}

\begin{convobox}{12}
    \noindent
    \jpkana{はい。}{Hai.} 
    \jpword{きんぎょ}{金魚}{Kingyo} \jpkana{を}{o} 
    \jpword{か}{飼}{ka}\jpkana{っていました。}{tte imashita.}

    \vspace{1.5em}
    \noindent{\Large \color{articolor} \textbf{Arti:} Ya. Aku memelihara ikan mas koki.}
\end{convobox}

\begin{convobox}{13}
    \noindent
    \jpword{こ}{子}{ko}\jpkana{どものころ、}{domo no koro,} 
    \jpword{す}{好}{su}\jpkana{きだった}{ki datta} 
    \jpkana{アニメ}{anime} \jpkana{は}{wa} 
    \jpword{なん}{何}{nan} \jpkana{でしたか？}{deshita ka?}

    \vspace{1.5em}
    \noindent{\Large \color{articolor} \textbf{Arti:} Waktu kecil, anime apa yang kamu sukai?}
\end{convobox}

\begin{convobox}{14}
    \noindent
    \jpkana{「キャプテン}{"Kyaputen} \jpword{つばさ}{翼}{Tsubasa}\jpkana{」という}{" to iu} 
    \jpkana{アニメ}{anime} \jpkana{が}{ga} 
    \jpword{す}{好}{su}\jpkana{きでした。}{ki deshita.}

    \vspace{1.5em}
    \noindent{\Large \color{articolor} \textbf{Arti:} Aku dulu suka anime bernama 'Captain Tsubasa'.}
\end{convobox}

\begin{convobox}{15}
    \noindent
    \jpword{こ}{子}{ko}\jpkana{どものころ、}{domo no koro,} 
    \jpword{にがて}{苦手}{nigate} \jpkana{だった}{datta} 
    \jpword{た}{食}{ta}\jpkana{べ}{be}\jpword{もの}{物}{mono} \jpkana{は}{wa} 
    \jpkana{ありましたか？}{arimashita ka?}

    \vspace{1.5em}
    \noindent{\Large \color{articolor} \textbf{Arti:} Waktu kecil, apakah ada makanan yang tidak kamu sukai?}
\end{convobox}

\begin{convobox}{16}
    \noindent
    \jpkana{ピーマン}{Piiman} \jpkana{と}{to} 
    \jpkana{にんじん}{ninjin} \jpkana{が}{ga} 
    \jpword{にがて}{苦手}{nigate} \jpkana{でした。}{deshita.}

    \vspace{1.5em}
    \noindent{\Large \color{articolor} \textbf{Arti:} Dulu aku tidak suka paprika hijau dan wortel.}
\end{convobox}

\begin{convobox}{17}
    \noindent
    \jpword{こ}{子}{ko}\jpkana{どものころ、}{domo no koro,} 
    \jpword{だれ}{誰}{dare} \jpkana{と}{to} 
    \jpkana{よく}{yoku} 
    \jpword{あそ}{遊}{aso}\jpkana{んでいましたか？}{nde imashita ka?}

    \vspace{1.5em}
    \noindent{\Large \color{articolor} \textbf{Arti:} Waktu kecil, sering bermain dengan siapa?}
\end{convobox}

\begin{convobox}{18}
    \noindent
    \jpword{きんじょ}{近所}{Kinjo} \jpkana{に}{ni} 
    \jpword{す}{住}{su}\jpkana{んでいた、}{nde ita,} 
    \jpkana{たいやきちゃんと}{Taiyaki-chan to} 
    \jpkana{よく}{yoku} 
    \jpword{あそ}{遊}{aso}\jpkana{んでいました。}{nde imashita.}

    \vspace{1.5em}
    \noindent{\Large \color{articolor} \textbf{Arti:} Aku dulu sering bermain dengan Taiyaki-chan yang tinggal di lingkungan sekitar.}
\end{convobox}

% --- LESSON BUFFER 3 ---
\begin{lessonbox}{Catatan Belajar: Mengubah Sifat Menjadi Masa Lalu!}
    Tahukah kamu? Kata sifat dan kata benda dalam bahasa Jepang akan berubah bentuk kalau peristiwanya sudah lewat lho!
    \begin{itemize}
        \item \textbf{Bentuk Sekarang:} \textit{Suki desu} (Suka) $\rightarrow$ \textbf{Bentuk Lampau:} \textit{Suki \textbf{deshita}} (Dulu suka).
        \item \textbf{Bentuk Sekarang:} \textit{Nigate desu} (Tidak suka/lemah) $\rightarrow$ \textbf{Bentuk Lampau:} \textit{Nigate \textbf{deshita}} (Dulu tidak suka).
    \end{itemize}
    Jadi ingat ya, kalau kamu mau cerita kesukaanmu pas masih TK atau SD, pakailah kata \textbf{〜でした (〜deshita)} di belakangnya!
\end{lessonbox}

% ================= PERCAKAPAN 19-26 =================
\begin{convobox}{19}
    \noindent
    \jpword{がっこう}{学校}{Gakkou} \jpkana{に}{ni} 
    \jpword{い}{行}{i}\jpkana{くのは}{ku no wa} 
    \jpword{す}{好}{su}\jpkana{きでしたか？}{ki deshita ka?} 
    \jpkana{きらいでしたか？}{Kirai deshita ka?}

    \vspace{1.5em}
    \noindent{\Large \color{articolor} \textbf{Arti:} Apakah kamu suka pergi ke sekolah? Atau benci?}
\end{convobox}

\begin{convobox}{20}
    \noindent
    \jpword{す}{好}{su}\jpkana{きでしたが、}{ki deshita ga,} 
    \jpword{がっこう}{学校}{gakkou} \jpkana{に}{ni} 
    \jpword{い}{行}{i}\jpkana{きたくない}{kitakunai} 
    \jpword{ひ}{日}{hi} \jpkana{も}{mo} 
    \jpkana{ありました。}{arimashita.}

    \vspace{1.5em}
    \noindent{\Large \color{articolor} \textbf{Arti:} Aku suka, tapi ada juga hari-hari di mana aku tidak ingin pergi ke sekolah.}
\end{convobox}

\begin{convobox}{21}
    \noindent
    \jpword{べんきょう}{勉強}{Benkyou} \jpkana{は}{wa} 
    \jpword{とくい}{得意}{tokui} \jpkana{でしたか？}{deshita ka?} 
    \jpword{にがて}{苦手}{nigate} \jpkana{でしたか？}{deshita ka?}

    \vspace{1.5em}
    \noindent{\Large \color{articolor} \textbf{Arti:} Apakah kamu jago belajar? Atau tidak bisa (lemah)?}
\end{convobox}

\begin{convobox}{22}
    \noindent
    \jpkana{ちょっと}{Chotto} 
    \jpword{にがて}{苦手}{nigate} \jpkana{でした。}{deshita.}

    \vspace{1.5em}
    \noindent{\Large \color{articolor} \textbf{Arti:} Sedikit tidak bisa (kurang pandai).}
\end{convobox}

\begin{convobox}{23}
    \noindent
    \jpword{がっこう}{学校}{Gakkou} \jpkana{には}{ni wa} 
    \jpword{なん}{何}{nan} \jpkana{で}{de} 
    \jpword{かよ}{通}{kayo}\jpkana{っていましたか？}{tte imashita ka?}

    \vspace{1.5em}
    \noindent{\Large \color{articolor} \textbf{Arti:} Pergi ke sekolah menggunakan (kendaraan) apa?}
\end{convobox}

\begin{convobox}{24}
    \noindent
    \jpword{たと}{例}{Tato}\jpkana{えば、}{eba,} 
    \jpword{くるま}{車}{kuruma} \jpkana{とか、}{toka,} 
    \jpkana{バス}{basu} \jpkana{とか。}{toka.}

    \vspace{1.5em}
    \noindent{\Large \color{articolor} \textbf{Arti:} Misalnya, dengan mobil, atau bus.}
\end{convobox}

\begin{convobox}{25}
    \noindent
    \jpword{しょうがっこう}{小学校}{Shougakkou} \jpkana{には、}{ni wa,} 
    \jpword{ある}{歩}{aru}\jpkana{いて}{ite} 
    \jpword{かよ}{通}{kayo}\jpkana{っていました。}{tte imashita.}

    \vspace{1.5em}
    \noindent{\Large \color{articolor} \textbf{Arti:} Ke SD, aku berjalan kaki.}
\end{convobox}

\begin{convobox}{26}
    \noindent
    \jpword{ちゅうがっこう}{中学校}{Chuugakkou} \jpkana{には、}{ni wa,} 
    \jpword{じてんしゃ}{自転車}{jitensha} \jpkana{で}{de} 
    \jpword{かよ}{通}{kayo}\jpkana{っていました。}{tte imashita.}

    \vspace{1.5em}
    \noindent{\Large \color{articolor} \textbf{Arti:} Ke SMP, aku naik sepeda.}
\end{convobox}

% --- LESSON BUFFER 4 ---
\begin{lessonbox}{Catatan Belajar: Naik Kendaraan vs Jalan Kaki}
    Coba lihat kalimat nomor 23. Ada grammar partikel \textbf{で (De)}!\\
    Dalam bahasa Jepang, kalau kita mau bilang pergi menggunakan alat transportasi, kita pakai partikel "De". 
    \begin{itemize}
        \item \textit{Kuruma \textbf{de}} (Menggunakan mobil)
        \item \textit{Jitensha \textbf{de}} (Menggunakan sepeda)
    \end{itemize}
    Tapi hati-hati! Kalau kamu pergi **dengan berjalan kaki**, jangan pakai "Ashi de" (dengan kaki) ya! Cukup bilang \textbf{歩いて (Aruite)} yang artinya "berjalan kaki" tanpa perlu partikel apa pun! Hebat kan?
\end{lessonbox}

% ================= PERCAKAPAN 27-36 =================
\begin{convobox}{27}
    \noindent
    \jpword{こ}{子}{ko}\jpkana{どものころ、}{domo no koro,} 
    \jpword{こわ}{怖}{kowa}\jpkana{かった}{katta} 
    \jpword{もの}{物}{mono} \jpkana{は}{wa} 
    \jpkana{ありましたか？}{arimashita ka?}

    \vspace{1.5em}
    \noindent{\Large \color{articolor} \textbf{Arti:} Waktu kecil, apakah ada hal yang menakutkan (buatmu)?}
\end{convobox}

\begin{convobox}{28}
    \noindent
    \jpkana{おばけ}{Obake} \jpkana{と}{to} 
    \jpword{かみなり}{雷}{kaminari} \jpkana{が}{ga} 
    \jpword{こわ}{怖}{kowa}\jpkana{かったです。}{katta desu.}

    \vspace{1.5em}
    \noindent{\Large \color{articolor} \textbf{Arti:} Dulu aku takut hantu dan petir.}
\end{convobox}

\begin{convobox}{29}
    \noindent
    \jpword{こ}{子}{ko}\jpkana{どものころの}{domo no koro no} 
    \jpword{ゆめ}{夢}{yume} \jpkana{は}{wa} 
    \jpword{なん}{何}{nan} \jpkana{でしたか？}{deshita ka?}

    \vspace{1.5em}
    \noindent{\Large \color{articolor} \textbf{Arti:} Waktu kecil, apa mimpimu (cita-cita)?}
\end{convobox}

\begin{convobox}{30}
    \noindent
    \jpword{うちゅうひこうし}{宇宙飛行士}{Uchuuhikoushi} \jpkana{でした。}{deshita.}

    \vspace{1.5em}
    \noindent{\Large \color{articolor} \textbf{Arti:} Dulu cita-citaku astronot.}
\end{convobox}

\begin{convobox}{31}
    \noindent
    \jpword{こ}{子}{ko}\jpkana{どものころの}{domo no koro no} 
    \jpword{じぶん}{自分}{jibun} \jpkana{に、}{ni,} 
    \jpkana{何か}{nani ka} 
    \jpword{い}{言}{i}\jpkana{いたいことは}{itai koto wa} 
    \jpkana{ありますか？}{arimasu ka?}

    \vspace{1.5em}
    \noindent{\Large \color{articolor} \textbf{Arti:} Apakah ada sesuatu yang ingin dikatakan kepada dirimu di masa kecil?}
\end{convobox}

\begin{convobox}{32}
    \noindent
    \jpword{ふとん}{布団}{Futon} \jpkana{の}{no} 
    \jpword{なか}{中}{naka} \jpkana{で}{de} 
    \jpkana{まんが}{manga} \jpkana{を}{o} 
    \jpword{よ}{読}{yo}\jpkana{むと}{mu to} 
    \jpword{め}{目}{me} \jpkana{が}{ga} 
    \jpword{わる}{悪}{waru}\jpkana{くなるよ、}{ku naru yo,} 
    \jpkana{と}{to} \jpword{い}{言}{i}\jpkana{いたいです。}{itai desu.}

    \vspace{1.5em}
    \noindent{\Large \color{articolor} \textbf{Arti:} Aku ingin bilang, "Kalau baca komik di dalam selimut, matamu akan rusak lho".}
\end{convobox}

\begin{convobox}{33}
    \noindent
    \jpword{きょう}{今日}{Kyou} \jpkana{は}{wa} 
    \jpkana{いろいろ}{iroiro} 
    \jpword{はな}{話}{hana}\jpkana{してくれて}{shite kurete} 
    \jpkana{ありがとうございました。}{arigatou gozaimashita.}

    \vspace{1.5em}
    \noindent{\Large \color{articolor} \textbf{Arti:} Terima kasih sudah menceritakan banyak hal hari ini.}
\end{convobox}

\begin{convobox}{34}
    \noindent
    \jpkana{たまには}{Tama ni wa} 
    \jpword{むかし}{昔}{mukashi} \jpkana{の}{no} 
    \jpword{はなし}{話}{hanashi} \jpkana{をするのも}{o suru no mo} 
    \jpkana{いいですね。}{ii desu ne.}

    \vspace{1.5em}
    \noindent{\Large \color{articolor} \textbf{Arti:} Sesekali membicarakan masa lalu itu menyenangkan ya.}
\end{convobox}

\begin{convobox}{35}
    \noindent
    \jpkana{また}{Mata} 
    \jpword{いっしょ}{一緒}{issho} \jpkana{に}{ni} 
    \jpword{はな}{話}{hana}\jpkana{しましょう。}{shimashou.}

    \vspace{1.5em}
    \noindent{\Large \color{articolor} \textbf{Arti:} Mari mengobrol bersama lagi nanti.}
\end{convobox}

\begin{convobox}{36}
    \noindent
    \jpkana{それでは、}{Sore dewa,} 
    \jpkana{また。}{mata.}

    \vspace{1.5em}
    \noindent{\Large \color{articolor} \textbf{Arti:} Kalau begitu, sampai jumpa.}
\end{convobox}

% --- SUMMARY BOX ---
\vspace{2em}
\begin{tcolorbox}[
    colback=blue!5!white,
    colframe=blue!80!black,
    boxrule=3pt,
    arc=5mm,
    title={\huge\bfseries \sffamily 🌟 Rangkuman Bab 5 🌟},
    fonttitle=\color{white},
    attach boxed title to top center={yshift=-4mm},
    boxed title style={colback=blue!80!black, rounded corners},
    left=5mm, right=5mm, top=7mm, bottom=5mm
]
\Large \textbf{Luar biasa! Kini kamu mahir bercerita masa lalu!} Yuk ingat grammar utamanya:
\begin{itemize}
    \item \textbf{Kata Kunci Masa Lalu:} \textit{Kodomo no koro} (Waktu kecil) / \textit{Mukashi} (Zaman dulu).
    \item \textbf{Akhiran Lampau:} Ubah \textit{Desu} menjadi \textbf{Deshita}. (Cth: \textit{Suki deshita} = Dulu suka).
    \item \textbf{Lebih suka yang mana?} Pakai grammar \textit{Dotchi no hou ga\dots} dan jawab dengan \textit{\dots hou ga suki deshita}.
    \item \textbf{Alat Transportasi:} Pakai partikel \textbf{De} (\textit{Basu de, Jitensha de}), kecuali jalan kaki gunakan \textbf{Aruite}.
\end{itemize}
Jangan lupa nasihat dari cerita di atas ya, jangan baca komik dalam gelap nanti matamu sakit! Sampai jumpa di bab selanjutnya yang tidak kalah seru!
\end{tcolorbox}
\newpage